\documentclass[12pt,a4paper,oneside]{report}
\usepackage[table]{xcolor}
% ===== PACKAGES =====
\usepackage[T1]{fontenc}
\usepackage[utf8]{inputenc}
\usepackage{mathptmx}           % Times New Roman font
\usepackage{anyfontsize}
\usepackage[english]{babel}

% Page layout: Left 1.25in, Right/Top/Bottom 1in
\usepackage[left=1.25in, right=1in, top=1in, bottom=1in, headheight=15pt]{geometry}

\usepackage{setspace}
\onehalfspacing                 % 1.5 line spacing

\usepackage{graphicx}
\usepackage{xcolor}
\usepackage{longtable,booktabs,array,multirow}
\usepackage{calc}
\usepackage{etoolbox}
\usepackage{float}
\usepackage{caption}
\usepackage{tocloft}
\usepackage{titlesec}
\usepackage{fancyhdr}
\usepackage{enumitem}
\usepackage{hyperref}
\usepackage{listings}

\hypersetup{hidelinks, pdfcreator={LaTeX}}

% ===== HEADINGS FORMAT =====

% First-level (Chapter): Times New Roman, 16 Bold, Left
% Leaves 2 blank lines (approx 36pt) between heading and first line of text
\titleformat{\chapter}[display]
  {\normalfont\fontsize{16}{19}\bfseries\raggedright}
  {Chapter \thechapter}
  {0pt}
  {\fontsize{16}{19}\bfseries\raggedright}
\titlespacing*{\chapter}{0pt}{-20pt}{36pt}  

% Second-level (Section): Times New Roman, 14 Bold, Left
% Leaves 2 blank lines before (36pt) and 1 blank line after (18pt)
\titleformat{\section}
  {\normalfont\fontsize{14}{17}\bfseries\raggedright}
  {\thesection}{1em}{}
\titlespacing*{\section}{0pt}{36pt}{18pt} 

% Third-level (Subsection): Times New Roman, 12 Bold, Left
% Leaves 2 blank lines before (36pt) and 1 blank line after (18pt)
\titleformat{\subsection}
  {\normalfont\fontsize{12}{15}\bfseries\raggedright}
  {\thesubsection}{1em}{}
\titlespacing*{\subsection}{0pt}{36pt}{18pt}

% Fourth level
\titleformat{\subsubsection}
  {\normalfont\fontsize{12}{15}\bfseries\raggedright}
  {\thesubsubsection}{1em}{}
\titlespacing*{\subsubsection}{0pt}{24pt}{12pt}

% ===== HEADER / FOOTER =====
\pagestyle{fancy}
\fancyhf{}
\renewcommand{\headrulewidth}{0.4pt}
\fancyhead[L]{\nouppercase{\leftmark}}   % Chapter title left
\fancyhead[R]{\thepage}                   % Page number right
\renewcommand{\chaptermark}[1]{\markboth{Chapter \thechapter\ - #1}{}}

% Plain style for chapter first page to enforce header rules universally
\fancypagestyle{plain}{%
  \fancyhf{}
  \fancyhead[L]{\nouppercase{\leftmark}}
  \fancyhead[R]{\thepage}
  \renewcommand{\headrulewidth}{0.4pt}
}

% Normal text justified (default in LaTeX report class)
\setlength{\parskip}{0pt}
\setlength{\parindent}{0pt}

% Caption formatting strictly by guidelines
\captionsetup[figure]{labelfont=bf, labelsep=colon, justification=centering, position=below}
\captionsetup[table]{labelfont=bf, labelsep=colon, justification=centering, position=above}

% TOC formatting (dotted leaders, nice numbering)
\renewcommand{\cftchapdotsep}{\cftdotsep}
\setcounter{tocdepth}{3}
\setcounter{secnumdepth}{3}

% Listings style
\lstset{
  basicstyle=\ttfamily\small,
  breaklines=true,
  frame=single,
  numbers=left,
  numberstyle=\tiny
}

\begin{document}

% ===== FRONT MATTER (Roman numerals) =====
\pagenumbering{roman}
\pagestyle{plain}

% ---------- TITLE PAGE ----------
\begin{titlepage}
\centering
\vspace*{0.5in}

{\fontsize{28}{34}\selectfont\bfseries SIVO\par}
\vspace{0.4in}

\includegraphics[width=1.86in]{vertopal_0a0373cd9b034ade9a141c927438222a/media/image1.jpeg}\par
\vspace{0.4in}

{\fontsize{14}{17}\selectfont\bfseries Muhammad Abubakr \quad BSE223219\par}
\vspace{0.15in}
{\fontsize{14}{17}\selectfont\bfseries Ahmad Abdullah Hashmi \quad BSE223105\par}
\vspace{0.15in}
{\fontsize{14}{17}\selectfont\bfseries Saail Abbas \quad BSE223144\par}
\vspace{0.6in}

{\fontsize{14}{17}\selectfont\bfseries Supervised By\par}
\vspace{0.1in}
{\fontsize{14}{17}\selectfont\bfseries Miss Khizra Bibi\par}
\vspace{0.5in}

{\fontsize{14}{17}\selectfont\bfseries Fall 2025\par}
\vspace{0.2in}
{\fontsize{14}{17}\selectfont\bfseries BSSE\par}
\vspace{0.4in}

{\fontsize{14}{17}\selectfont\bfseries Department of Software Engineering\par}
\vspace{0.15in}
{\fontsize{14}{17}\selectfont\bfseries Capital University of Science \& Technology, Islamabad\par}

\vfill
\end{titlepage}

% ---------- PROJECT REPORT PAGE ----------
\newpage
\thispagestyle{empty}

\vspace*{0.4in}

% --- Header with Centered Title and Right-Aligned Logo ---
\noindent
\begin{minipage}{0.33\textwidth}
\strut
\end{minipage}%
\begin{minipage}{0.33\textwidth}
\centering
{\fontsize{20}{24}\selectfont\bfseries Project Report\par}
\end{minipage}%
\begin{minipage}{0.33\textwidth}
\raggedleft
\includegraphics[width=1.22in]{vertopal_0a0373cd9b034ade9a141c927438222a/media/image2.jpg}
\end{minipage}

\vspace{0.8in}

% --- Version and Members Row ---
\noindent
\begin{tabular}{|>{\columncolor{black}\color{white}\bfseries}p{1.2in}|p{1.5in}|}
\hline
VERSION & V X.0 \\
\hline
\end{tabular}
\hfill
\begin{tabular}{|>{\columncolor{black}\color{white}\bfseries\centering\arraybackslash}p{1.4in}|p{0.6in}|}
\hline
NUMBER OF \newline MEMBERS & 3 \\
\hline
\end{tabular}

\vspace{0.4in}

% --- Title Row ---
\noindent
\begin{tabular}{|>{\columncolor{black}\color{white}\bfseries}p{1.2in}|p{4.72in}|}
\hline
TITLE & SIVO: Two-Way Sign Language Translator \\
\hline
\end{tabular}

\vspace{0.4in}

% --- Supervisor Row ---
\noindent
\begin{tabular}{|>{\columncolor{black}\color{white}\bfseries}p{1.7in}|p{4.22in}|}
\hline
SUPERVISOR NAME & Miss Khizra Bibi \\
\hline
\end{tabular}

\vspace{0.6in}

% --- Members Table ---
\noindent
\begin{tabular}{|p{2in}|p{1.3in}|p{2.3in}|}
\hline
\rowcolor{black}\textcolor{white}{\textbf{MEMBER NAME}} & \textcolor{white}{\textbf{REG. NO.}} & \textcolor{white}{\textbf{EMAIL ADDRESS}} \\
\hline
Muhammad Abubakr & BSE223219 & \\ \hline
Ahmad Abdullah Hashmi & BSE223105 & \\ \hline
Saail Abbas & BSE223144 & \\ \hline
\end{tabular}

\vspace{0.6in}

% --- Signatures Section ---
\noindent
\begin{minipage}[t]{0.5\textwidth}
    % Black Header Box for Members
    \begin{tabular}{|>{\columncolor{black}\color{white}\bfseries\centering\arraybackslash}p{2.5in}|}
    \hline
    MEMBERS' SIGNATURES \\
    \hline
    \end{tabular}
    
    \vspace{0.6in}
    \rule{2.5in}{0.8pt} % Signature Line 1
    
    \vspace{0.4in}
    \rule{2.5in}{0.8pt} % Signature Line 2
    
    \vspace{0.4in}
    \rule{2.5in}{0.8pt} % Signature Line 3
\end{minipage}%
\hfill
\begin{minipage}[t]{0.45\textwidth}
    \vspace{0.2in} % Adjust vertical alignment to match the left side
    \centering
    % Framed Box for Supervisor
    \framebox{\parbox[c][1.8in][c]{2.5in}{\centering \textbf{Supervisor's Signature}}}
\end{minipage}

% ---------- APPROVAL CERTIFICATE ----------
\chapter*{Approval Certificate}
\addcontentsline{toc}{chapter}{Approval Certificate}
\markboth{Approval Certificate}{}

This project, entitled as ``SIVO: Two-Way Sign Language Translator'' has been approved for the award of

\vspace{0.3in}
\begin{center}
{\fontsize{14}{17}\selectfont\bfseries Bachelors of Science in Software Engineering\par}
\end{center}

\vspace{0.4in}
\textbf{Committee Signatures:}

\vspace{0.4in}
Supervisor: \rule{3in}{0.4pt}

\hspace{0.5in}(Supervisor Name Dr / Mr / Ms)

\vspace{0.4in}
Project Coordinator: \rule{3in}{0.4pt}

\hspace{0.5in}(Mr. Ibrar Arshad)

\vspace{0.4in}
Head of Department: \rule{3in}{0.4pt}

\hspace{0.5in}(Dr. Nadeem Anjum)

\vspace{0.5in}
\textbf{MEMBERS' SIGNATURES}

% ---------- DECLARATION ----------
\newpage
\thispagestyle{plain} % Ensure page number is at the bottom center, or empty if preferred

\chapter*{Declaration}
\addcontentsline{toc}{chapter}{Declaration}
\markboth{Declaration}{}

\vspace{0.2in}

\noindent We, hereby, declare that ``No portion of the work referred to, in this project has been submitted in support of an application for another degree or qualification of this or any other university/institute or other institution of learning''. It is further declared that this undergraduate project, neither as a whole nor as a part thereof has been copied out from any sources, wherever references have been provided.

\vspace{2.5in}

\begin{flushright}
\begin{minipage}{0.45\textwidth}
    \begin{tabular}{|>{\columncolor{black}\color{white}\bfseries\centering\arraybackslash}p{2.5in}|}
    \hline
    MEMBERS' SIGNATURES \\
    \hline
    \end{tabular}
    
    \vspace{0.6in}
    \rule{2.5in}{0.8pt}
    
    \vspace{0.4in}
    \rule{2.5in}{0.8pt}
    
    \vspace{0.4in}
    \rule{2.5in}{0.8pt}
\end{minipage}
\end{flushright}
% ---------- ACKNOWLEDGEMENTS ----------
\chapter*{Acknowledgements}
\addcontentsline{toc}{chapter}{Acknowledgements}
\markboth{Acknowledgements}{}

We would like to express our sincere gratitude to all those who contributed to the successful completion of this project.

First and foremost, we thank our project supervisor for providing invaluable guidance, support, and encouragement throughout the development of SIVO.

We are also grateful to our faculty members and the Department of Computer Science for providing us with the necessary resources and facilities.

Special thanks to the deaf community members who participated in our testing and provided valuable feedback that helped improve the application.

Finally, we thank our families and friends for their continuous support and understanding during this journey.

% ---------- DEDICATION ----------
\chapter*{Dedication}
\addcontentsline{toc}{chapter}{Dedication}
\markboth{Dedication}{}

\textit{This project is dedicated to the deaf and mute community of Pakistan, with the hope that SIVO will help bridge the communication gap and create a more inclusive society.}

% ---------- EXECUTIVE SUMMARY ----------
\chapter*{Executive Summary}
\addcontentsline{toc}{chapter}{Executive Summary}
\markboth{Executive Summary}{}

SIVO is a mobile application designed to facilitate real-time, two-way communication between deaf and hearing individuals in Pakistan. The application addresses a significant gap in the market by providing the first dedicated Pakistani Sign Language (PSL) translator.

The system operates in two modes: Sign-to-Speech, which converts hand gestures captured through the mobile camera into spoken audio using artificial intelligence, and Speech-to-Sign, which transforms spoken words into animated sign language performed by a 3D avatar.

Key features of SIVO include real-time translation, grammar correction using the Grok API, conversation history management, and a freemium monetization model. The application is built using React Native for the mobile frontend and Python Flask for the backend AI processing.

The primary target users are business owners and employees in organizations who need to communicate with deaf colleagues or customers. By eliminating communication barriers, SIVO aims to promote workplace inclusion and improve professional interactions for the deaf community.

This report documents the complete development lifecycle of SIVO, including requirements analysis, system design, implementation, and deployment procedures.

% ---------- TABLE OF CONTENTS ----------
\newpage
\tableofcontents

% ---------- LIST OF FIGURES ----------
\newpage
\addcontentsline{toc}{chapter}{List of Figures}
\listoffigures

% ---------- LIST OF TABLES ----------
\addcontentsline{toc}{chapter}{List of Tables}
\listoftables

% ===== MAIN MATTER (Arabic numerals) =====
\newpage
\pagenumbering{arabic}
\setcounter{page}{1}
\pagestyle{fancy}

% ==========================================================
% CHAPTER 1: INTRODUCTION
% ==========================================================
\chapter{Introduction}

In Pakistan, many people are deaf and silent, and they face many challenges that limit their ability to communicate effectively. First, in the environment of business and workplace, deaf persons struggle to engage in meetings, distribute presentations, or serve customers efficiently, which affects their professional development and confidence. Secondly, listeners also find it difficult to interact naturally with their deaf colleagues or customers, leading to social and professional isolation for both sides. Third, there is still no reliable and accessible digital solution that can translate specifically PSL sign language into speech in real time. There are several applications and websites for languages like ASL and BSL but not for local PSL.

The project, called SIVO, is designed to solve this problem. SIVO is a mobile application that works like Google Translator, but instead of converting a spoken language into another, it turns Pakistani Sign Language (PSL) into speech, and turns speech back into PSL. For example, when a deaf person signs in front of the camera, the application will turn it into clear Urdu or English speech. And when a hearing person speaks, the application will turn it into PSL gestures on the screen. In this way, both sides can eventually understand each other without the requirement of a third person.

To make it possible, a dataset of normal PSL signals will first be collected. Then, a machine learning model will be trained to correctly identify these signs. After that, a mobile application will be created that adds camera to sign input and uses speech technology for voice translation. Finally, the application will be tested with business people and employees to ensure that it works in real situations. Step by step, SIVO will transform from an idea into an actual tool that helps thousands of people to communicate better in their jobs, businesses, and everyday life.

\section{Project Introduction}

“The goal of this project was to develop SIVO, a mobile application capable of bidirectional translation between Pakistani Sign Language (PSL) and spoken language”.

To build this system, a diverse PSL dataset was collected, incorporating pre-existing legacy recordings, user webcam captures, and 1080p mobile videos across more than 140 word classes. Rather than training on raw images using Convolutional Neural Networks (CNNs), the system utilized MediaPipe Holistic to extract dynamic facial, pose, and hand landmarks. This sequential data was preprocessed using a sliding window technique and extensive data augmentation to train a robust GRU-based Recurrent Neural Network capable of recognizing temporal, continuous signs. The final model was optimized using TensorFlow Lite (TFLite) for high-speed inference.


Following word prediction, a custom "Smart Match" algorithm was implemented to map recognized words into grammatically appropriate sentences based on a conversational knowledge base (with third-party APIs available as a fallback). For the reverse communication flow (Speech-to-Sign), on-device speech recognition transcribes spoken language. An NLP pipeline lemmatizes the text and maps the resulting words to a bundled dictionary of pre-recorded 3D avatar videos, displaying them as a seamless visual sequence.
All of these phases—data collection, model training, prediction, sentence formation, and speech conversion—have been successfully integrated into a real-time, cross-platform mobile application. 


A deaf user can sign in front of the camera, and a hearing user can simply speak, enabling smooth two-way communication. The primary beneficiaries of this project are businesses and organizations in Pakistan, where communication barriers with deaf individuals have historically hindered productivity and workplace inclusion \hyperlink{ref1}{[1]}, \hyperlink{ref7}{[7]}.



\textbf{Business People:} SIVO helps business owners and entrepreneurs communicate smoothly with clients, employees, and partners. It removes the communication barrier during meetings, customer interactions, and daily operations, making business activities more inclusive and efficient.

\textbf{Employees in Large Organizations:} SIVO supports employees, especially those in corporate or office environments, by enabling clear two-way communication between deaf and hearing staff. It ensures better teamwork, understanding in meetings, and overall workplace accessibility.

\subsection{System Design Document}

\textbf{a) App Concept}

The main concept is to create a two-way, real-time communication bridge.

\begin{itemize}
  \item \textbf{Sign-to-Speech:} The system uses machine learning, specifically MediaPipe and feed-forward network (MLP) model, to analyze the user's hand, currency, and facial gestures through a live video feed, converting them into Urdu or English.
  \item \textbf{Speech-to-Sign:} To display the correct PSL signals on the screen, a custom animation/gesture engine uses speech-to-text processing, which translates the spoken word into visual communication.
\end{itemize}

The creation of the application includes PSL gestures, advanced AI model training, and a simple, user-friendly mobile interface with comprehensive data collection.

\textbf{b) Target Audience}

\begin{itemize}
  \item \textbf{Primary:} Business owners and employees in big organizations who are deaf and mute individuals who use Pakistani Sign Language (PSL).
  \item \textbf{Secondary:} Family members, teachers, businesses, healthcare professionals, and anyone who regularly interacts with the deaf community.
\end{itemize}

\textbf{c) App Flow (User Journey)}

The following figure illustrates the application flow and user journey through the SIVO system. Figure~\ref{fig:appflow} demonstrates the complete user journey through the SIVO application, 

\begin{figure}[H]
\centering
\includegraphics[width=0.75\textwidth]{vertopal_0a0373cd9b034ade9a141c927438222a/media/image3.jpg}
\caption{SIVO Application Flow Diagram}
\label{fig:appflow}
\end{figure}

\begin{itemize}
  \item \textbf{Opening the App:} User lands on a clean, dual-mode screen (Sign to Speech or Speech to Sign).
  \item \textbf{Sign-to-Speech:}
  \begin{itemize}
    \item Deaf User taps the ``Sign'' button.
    \item Camera view opens.
    \item User performs their sign language.
    \item User taps ``Convert''.
    \item The application displays the translated text and immediately speaks the sentence in the selected output language (Urdu/English).
  \end{itemize}
  \item \textbf{Speech-to-Sign:}
  \begin{itemize}
    \item Hearing User taps the ``Speech'' button.
    \item User speaks into the microphone.
    \item The application instantly processes the speech and displays the corresponding sign language gestures as a smooth, animated sequence on the screen.
  \end{itemize}
\end{itemize}

\textbf{d) Project Scope}

The scope of this project is to design and develop SIVO, a mobile application that bridges the communication gap between the deaf/mute community and the speaking population. The application will provide real-time translation of Pakistani Sign Language (PSL) into spoken Urdu/English and convert spoken Urdu/English into PSL gestures using animations or visuals.

\textbf{Inclusions:}
\begin{itemize}
  \item Real-time sign recognition through the mobile camera.
  \item Conversion of PSL gestures into speech (Urdu and English) (300 to 400 words).
  \item Conversion of spoken Urdu/English into PSL animations.
  \item A simple, user-friendly interface accessible to all age groups.
  \item Support for mobile platforms (Android initially, iOS later if time allows).
\end{itemize}

\textbf{Exclusions:}
\begin{itemize}
  \item Support for international sign languages beyond PSL (e.g., ASL, BSL).
  \item Advanced personalization features (like saving conversations, offline training).
  \item No more than 400 words, but the dictionary can be increased with time.
\end{itemize}

\textbf{Boundaries:} This project will focus only on basic daily-life conversations and business-related communication between deaf/mute users and normal speakers within Pakistan. The solution will be scalable in the future to add more languages, advanced features, and cross-platform support.

\section{Existing Examples / Solutions}

Table~\ref{tab:existing} compares the existing sign language translation solutions available in the market. It highlights the languages supported, the key functionalities offered, and the shortcomings of each product. This comparison clearly shows the gap in the market regarding Pakistani Sign Language (PSL), which SIVO aims to fill.

\begin{longtable}{|p{1.0in}|p{1.1in}|p{1.7in}|p{1.7in}|}
\caption{Comparison of Existing Sign Language Solutions}\label{tab:existing}\\
\hline
\textbf{Product/System} & \textbf{Sign Language} & \textbf{Functionality} & \textbf{Shortcomings / Missing Feature} \\
\hline
\endfirsthead
\hline
\textbf{Product/System} & \textbf{Sign Language} & \textbf{Functionality} & \textbf{Shortcomings / Missing Feature} \\
\hline
\endhead
Sign-Speak (Website) & American Sign Language (ASL) & Two-way translation (Sign to Text/Speech) & Crucially, it does not support PSL, making it unusable for the core Pakistani market. \\
\hline
Hand Talk (App) & American Sign Language (ASL) & Often a dictionary or isolated word recognition. & Focus solely on ASL. Limited to isolated words, not continuous, natural sentences. \\
\hline
Lingvano (App) & BSL, ASL, Australian Sign Language & Subscription-based app with video lessons, interactive exercises, self-practice tools. & Paid App, free features are too limited. Problems with media not playing video and pictures during lesson. Progress lost when restarting. \\
\hline
ASL Bloom (App) & American Sign Language (ASL) only & 120 video lessons covering more than 1300 signs and sentences. Searchable dictionary of signs. Interactive Quizzes. & Premium Subscription Only. Follow sequence of videos in fixed order. For advanced ASL it fails. No Real Time Gestures only video lessons. \\
\hline
Let's Talk Sign (App) & Indian Sign Language (ISL) only & Real Time speech to Sign, Text to Sign, Text to speech, Camera to scan text, Ready-made sentences. & Scan option only scans printed text. Offline features are more simple. Real-time recognition can cause errors due to gestures and speed of speech. \\
\hline
\end{longtable}

The main drawback of existing solutions is the total lack of support for Pakistani Sign Language (PSL) in a real-time, continuous, two-way translation application. SIVO aims to fill this gap by creating the first dedicated, strong application for the PSL community. While proving products such as Sign-Speak is technically possible, the focus on PSL data collection and modeling is a unique selling point that works in a large-scale, neglected market.

\section{Business Scope}

SIVO has an important ability to be a viable, life-changing commercial product due to focusing on a large market.

\textbf{Target Market and Professional Capacity:}

\begin{itemize}
  \item \textbf{Core Users:} SIVO's primary users are business owners, salespersons, and workers in mid-size and large companies. These professionals often face serious challenges in connecting with deaf/silent customers or colleagues, causing communication intervals, productivity loss, and lost opportunities.
  \item \textbf{Professional Environment:} The application is particularly valuable in customer-supported roles such as retail counters, sales offices, service desks, and corporate workplaces, where direct communication is important for smooth operation and customer satisfaction.
\end{itemize}

Table~\ref{tab:leancanvas} presents the Lean Canvas Template for the SIVO project. It summarizes the core problem, proposed solution, unique value proposition, customer segments, revenue streams, and other essential business elements that define the project's commercial viability and market strategy.

\begin{longtable}{|p{1.6in}|p{4.3in}|}
\caption{Lean Canvas Template}\label{tab:leancanvas}\\
\hline
\textbf{Section} & \textbf{Description} \\
\hline
\endfirsthead
\hline
\textbf{Section} & \textbf{Description} \\
\hline
\endhead
Problem & 1. Deaf and mute individuals struggle to communicate in workplaces and business environments, especially with managers and customers. 2. Hearing people find it difficult to understand or respond to sign language, creating communication gaps and inefficiency. 3. There is no real-time, reliable mobile solution in Pakistan that can translate PSL sign language into speech and vice versa. \\
\hline
Solution & Develop a real-time mobile application that converts sign gestures into spoken words and speech into animated signs. Use custom-built dataset for training of PSL sign Language. \\
\hline
Unique Value Proposition (UVP) & A two-way communication app that performs sentence-level translation between Pakistani Sign Language (PSL) and speech, enabling instant, natural communication between deaf and hearing users. \\
\hline
Unfair Advantage & The custom dataset is created completely by the team with real-world sign variations. No one else can copy or replicate it easily. This gives a unique, proprietary edge in PSL-based communication systems. \\
\hline
Customer Segments & Deaf and mute employees in professional and business environments. Business owners, managers, and salespeople who work with or serve deaf clients. \\
\hline
Existing Alternatives & Sign-Speak (website) - Built for ASL, not PSL. SignAll / HandTalk / other apps - Mostly trained on ASL/other languages and often only do word-level or demo translations. \\
\hline
Key Metrics & Translation accuracy (both sign to speech and speech to sign). Response time during real-time translation. Number of active users and retention rate. User satisfaction and feedback ratings. \\
\hline
High-Level Concept & ``SIVO is like a Google Translator for Sign Language - but designed specifically for real-time, two-way workplace communication.'' \\
\hline
Channels & App stores (Google Play, Apple App Store). Social media campaigns and professional awareness events. \\
\hline
Early Adopters & Business people and employees within the team's own family and close network who face communication barriers while interacting with deaf coworkers or clients. \\
\hline
Cost Structure & Data collection and labeling. Server hosting and maintenance. \\
\hline
Revenue Structure & Freemium model with a paid premium version. In-app ads or accessibility-focused subscriptions. \\
\hline
\end{longtable}

\subsection{Monetization Model}

Since this project is based on a mobile application that translates Pakistani Sign Language (PSL) into text/speech, several ways can be explored to make it sustainable and profitable:

\begin{itemize}
  \item \textbf{Freemium (Main Model):} The application will be free to use for basic translations (daily conversation gestures). Users who need advanced features, like business vocabulary packs, offline usage, or premium voices, can upgrade to the paid version.
  \item \textbf{Advertising (Optional):} Light, non-intrusive ads can be shown in the free version (for example: small banners at the bottom). But since this application serves a sensitive audience, ads should never interrupt or annoy the user.
\end{itemize}

\section{Useful Tools and Technologies}

\textbf{1. Programming Language}
\begin{itemize}
  \item Python will be used for building the sign recognition model because it is the best choice for AI/ML projects. It has powerful libraries like TensorFlow, Keras, and OpenCV which make gesture recognition easier.
  \item For the mobile application itself, Kotlin will be used because it allows building Android applications.
\end{itemize}

\textbf{2. Development Environment}
\begin{itemize}
  \item PyCharm / Jupyter Notebook for building and testing the AI model (because they are great for debugging and experimenting with ML code).
  \item Android Studio / VS Code for Kotlin application development (because they provide strong support for mobile UI design and debugging).
\end{itemize}

\textbf{3. Database}
\begin{itemize}
  \item Firebase will be used because it is cloud-based, real-time, and integrates smoothly with Flutter applications. It will store user data, application preferences, and sign vocabulary sets.
\end{itemize}

\textbf{4. Operating Systems}
\begin{itemize}
  \item Android is the main target since the application will be mobile-first.
  \item Later, the project can also be extended to Web for schools and offices.
\end{itemize}

\textbf{5. Network Protocol}
\begin{itemize}
  \item HTTP/HTTPS will be used to connect the mobile application with cloud services (Firebase, APIs, or AI model updates).
  \item This ensures secure communication, especially when sending/receiving sign translation data.
\end{itemize}

\section{Project Work Breakdown}

Table~\ref{tab:wbs} presents the Project Work Breakdown Structure for the SIVO development. It divides the entire project into logical phases and sub-phases, assigns lead members and supporting members for each task, and provides a clear overview of responsibilities. This structured breakdown ensures efficient task distribution and accountability among team members.

\begin{longtable}{|p{1.2in}|p{2.2in}|p{1in}|p{1.3in}|}
\caption{Project Work Breakdown Structure}\label{tab:wbs}\\
\hline
\textbf{Phase \& Sub-Phase} & \textbf{Task Description} & \textbf{Lead Member} & \textbf{Supporting Members} \\
\hline
\endfirsthead
\hline
\textbf{Phase \& Sub-Phase} & \textbf{Task Description} & \textbf{Lead Member} & \textbf{Supporting Members} \\
\hline
\endhead
\multicolumn{4}{|l|}{\textbf{P1. Requirements \& Analysis}} \\
\hline
P1.1 & Finalize PSL Vocabulary & Saail & All \\
\hline
P1.2 & Define System Architecture (Database, Model Structure) & Ahmed & Abubakar \\
\hline
\multicolumn{4}{|l|}{\textbf{P2. Design \& Data Acquisition}} \\
\hline
P2.1 & Data Collection Tools \& Environment Setup & Ahmed & All \\
\hline
P2.2 & PSL Video Dataset Gathering \& Labeling & All & All \\
\hline
P2.3 & Raw Data Preprocessing Pipeline (MediaPipe) & Abubakar & Saail \\
\hline
P2.4 & Sign-To-Speech Logic Design & Abubakar & Saail \\
\hline
\multicolumn{4}{|l|}{\textbf{P3. Implementation}} \\
\hline
P3.1 & Mobile App Frontend \& UI/UX Development & Abubakar & Saail (Feedback) \\
\hline
P3.2 & Word-Level Model Training (MLP or best accuracy) & Abubakar & Ahmed \\
\hline
P3.3 & Sentence-Level Model (Grok API) & Saail & Ahmed \\
\hline
P3.4 & Grammar Correction Module Development & Abubakar & Saail \\
\hline
\multicolumn{4}{|l|}{\textbf{P4. Testing \& Integration}} \\
\hline
P4.1 & Model Optimization \& TFLite Conversion & Abubakar & Saail \\
\hline
P4.2 & Backend Integration (Sign to Speech Modules) & Abubakar & Saail \\
\hline
P4.3 & Alpha/Beta Testing and Quality Assurance (QA) & All & All \\
\hline
\multicolumn{4}{|l|}{\textbf{P5. Deployment \& Finalization}} \\
\hline
P5.1 & Final Bug Fixing and Performance Tuning & Abubakar & Saail \\
\hline
P5.2 & Documentation, Report \& Presentation & All & All \\
\hline
\end{longtable}

Figure~\ref{fig:sysdev} illustrates the overall SIVO system development architecture, showing the key components and their interactions within the application ecosystem. It provides a high-level visual summary of how the various development phases integrate to produce the final functional system.

\begin{figure}[H]
\centering
\includegraphics[width=0.75\textwidth]{vertopal_0a0373cd9b034ade9a141c927438222a/media/image4.jpg}
\caption{SIVO System Development Overview}
\label{fig:sysdev}
\end{figure}

\section{Project Timeline}

Table~\ref{tab:timeline} presents the detailed project timeline with phase IDs, components, allocated time in weeks, percentage distribution, and justification for the time allocation. This breakdown helps visualize the distribution of effort across the different stages of the project lifecycle.

\begin{longtable}{|p{0.6in}|p{1.2in}|p{0.7in}|p{0.8in}|p{2.1in}|}
\caption{Project Timeline}\label{tab:timeline}\\
\hline
\textbf{Phase ID} & \textbf{Component} & \textbf{Time (Weeks)} & \textbf{Percentage} & \textbf{Justification} \\
\hline
\endfirsthead
\hline
\textbf{Phase ID} & \textbf{Component} & \textbf{Time (Weeks)} & \textbf{Percentage} & \textbf{Justification} \\
\hline
\endhead
P1 & Requirement Gathering \& Analysis & 2 & 6\% & Defining the scope, PSL vocabulary, and final system architecture. \\
\hline
P2 & Design \& Data Acquisition & 8 & 30\% & Highest Allocation: Includes critical PSL Data Collection, the most time-consuming task. \\
\hline
P3 & Implementation \& Model Training & 13 & 31.25\% & Focused on model training and mobile development. \\
\hline
P4 & Testing \& Integration & 3 & 18.75\% & Integration of ML model into mobile app and comprehensive testing. \\
\hline
P5 & Deployment \& Finalization & 2 & 6.5\% & Final bug fixing, report writing, and project handover. \\
\hline
\textbf{Total} & --- & \textbf{28 Weeks} & \textbf{100\%} & Approximately 8 months project duration. \\
\hline
\end{longtable}

The following Gantt chart provides a visual representation of the project timeline. Figure~\ref{fig:gantt} illustrates the project timeline in Gantt chart format, showing the duration and overlap of various project phases from requirements gathering through deployment and finalization. This visualization helps in tracking progress and understanding dependencies between different phases.

\begin{figure}[H]
\centering
\includegraphics[width=0.8\textwidth]{vertopal_0a0373cd9b034ade9a141c927438222a/media/image5.jpg}
\caption{Project Gantt Chart}
\label{fig:gantt}
\end{figure}

% ==========================================================
% CHAPTER 2: REQUIREMENT SPECIFICATION AND ANALYSIS
% ==========================================================
\chapter{Requirement Specification and Analysis}

This section provides a comprehensive explanation of the complete functionality of the SIVO application. It encompasses the major features, user stories, test cases, and the operational workflow of the entire system including sign-to-speech conversion, speech-to-sign translation, grammar correction mechanisms, user interface design, performance optimization, and security implementations. Each epic represents a group of related features, while each user story describes a specific user requirement. Every story is accompanied by a test case to verify its correct functionality.

\section{SIVO Epics}

\subsection{Epic 1: User Account and Security}

\textbf{Description:}

